% Section II: Related Work (deep learning) -- multi-task learning, deep networks
% for Wi-Fi fingerprinting, inputs and regularization, trees vs. networks and
% hybrids, and the gap. What each paper did or reported is taken from
% docs/literature_review_phase2.md and docs/literature_matrix_phase2.csv
% (with source locations). No result of our own appears in this section.
\section{Related Work}\label{sec:related}

\emph{Multi-task learning.} Caruana~\cite{caruana1997multitask} showed that
related tasks learned in parallel with a shared representation can help each
other. In deep networks this is usually \emph{hard parameter sharing}: shared
hidden layers and one output layer per task~\cite{ruder2017overview}. It
reduces overfitting, but can fail when the tasks are not closely
related~\cite{ruder2017overview}, and the results depend strongly on the
weights of the task losses~\cite{kendall2018multi}.

\emph{Deep networks for fingerprinting.} Phase~1 reviewed the autoencoder
networks of~\cite{nowicki2017loweffort,kim2018scalable} and
CNNLoc~\cite{song2019cnnloc}. Kim's single-input multi-output (SIMO)
network~\cite{kim2018simo} has the design of our brief: a shared hidden layer, a
floor head and a coordinate regression head (on another dataset). On
UJIIndoorLoc, with its published split and averages over 10 runs, a SIMO-style
MLP reached a mean 2D error of 9.20~m, but a single-task position network
reached 8.25~m~\cite{laska2022multitask}; a position network started from a
trained building-and-floor network was slightly better than one network
trained jointly~\cite{li2025hierarchical}. Hierarchical recurrent networks
found 95.23\% of the floors, with a setting that appears to have been read off
the same test results~\cite{elesawi2021hierarchical}. The design closest to ours, a shared
trunk with three heads and a strongly weighted position
loss~\cite{li2026hybrid}, tuned part of its system on the file it reports as
its test set. Results are hard to compare: metrics differ (2D or 3D),
settings are often chosen on the test file, and one network gave 9.52~m on the
published split but 2.36~m on a random re-split~\cite{kargarbarzi2024edge}; a
review names the lack of a common evaluation framework as a main
problem~\cite{alhomayani2020deep}.

\emph{Inputs and regularization.} Scaling each reading to $[0,1]$ with not
heard as 0~\cite{kargarbarzi2024edge,li2026hybrid}, adding a 0/1 ``detected''
channel per WAP and hiding WAPs at random during training (a
preprint~\cite{eftekharizadeh2026physics}) have all been used before; hiding
inputs is a form of dropout~\cite{srivastava2014dropout}. We train with
ReLU~\cite{nair2010rectified,goodfellow2016deep}, AdamW~\cite{kingma2015adam,%
loshchilov2019decoupled} and dropout~\cite{srivastava2014dropout}.

\emph{Trees, networks and hybrids.} On medium-sized tabular benchmarks, tree
models remain state of the art and MLPs suffer from uninformative
features~\cite{grinsztajn2022tree}. That benchmark leaves out signal data like
ours and missing values, so it only suggests that a network may not beat our
Phase~1 Random Forest. Wide-and-deep models train a network jointly with
hand-made features~\cite{cheng2016wide}; on tabular data, an ensemble of
networks and XGBoost beat XGBoost alone~\cite{shwartzziv2022tabular}; and
fingerprinting hybrids match scans with weighted kNN in a learned
space~\cite{park2024wknn}, or average such a kNN position with the network's
own, with the weight chosen on the file reported as the test
set~\cite{li2026hybrid}.

\emph{Gap.} We found no study on UJIIndoorLoc that (i) compares a multi-task
network with tuned tree models and with single-task networks while tuning every
setting by user-group CV and using the test file only at the end, and (ii)
combines learned location profiles of the WAPs (our Phase~1 AP-affinity) and a
tree model with a multi-task network, with an ablation that measures their
value. Phase~2 does both.
